% Generated draft: requires mathematical and source review.
\documentclass[11pt]{article}
\usepackage[margin=25mm]{geometry}
\usepackage{fontspec}
\setmainfont{Latin Modern Roman}
\usepackage{amsmath,amssymb,mathtools,unicode-math}
\setmathfont{Latin Modern Math}
\usepackage{enumitem,needspace,xurl,hyperref}
\setlength{\parindent}{0pt}
\setlength{\parskip}{5pt}
\setlength{\emergencystretch}{3em}
\newcommand{\N}{\mathbb{N}}
\newcommand{\Z}{\mathbb{Z}}
\newcommand{\Q}{\mathbb{Q}}
\newcommand{\R}{\mathbb{R}}
\newcommand{\C}{\mathbb{C}}
\newcommand{\F}{\mathbb{F}}
\newcommand{\A}{\mathbb{A}}
\newcommand{\PP}{\mathbb{P}}
\newcommand{\E}{\mathbb{E}}
\newcommand{\eps}{\varepsilon}
\newcommand{\dd}{\,\mathrm{d}}
\begin{document}
\section*{UnsolvedMath ID: 30004882}
\textbf{Source problem code:} catalogue entry
\subsection{Definitions and mathematical statement}
The supplied question concerns a quiver \(Q\), its preprojective algebra \(\Pi_Q\), and its cohomological Hall algebra \(A_{\Pi_Q}\). The record does not specify the allowed quivers, coefficient field, or precise construction of this cohomological Hall algebra.

For clarity of notation, under the usual finite-quiver convention, write \(Q=(Q_0,Q_1,s,t)\), where \(Q_0\) and \(Q_1\) are the vertex and arrow sets and \(s,t\) give the source and target of each arrow. Its doubled quiver \(\overline Q\) is obtained by adjoining an arrow \(a^*:t(a)\to s(a)\) for every \(a\in Q_1\). Over a field \(k\), the ordinary preprojective algebra is
\[
\Pi_Q=k\overline Q\Big/\left\langle\sum_{a\in Q_1}(aa^*-a^*a)\right\rangle,
\]
where \(k\overline Q\) is the path algebra. This conditional definition explains standard notation; it does not resolve whether this finite-quiver convention, this coefficient field, or this ordinary version is intended by the source.

\textbf{Problem.} Determine the structure of the entire cohomological Hall algebra \(A_{\Pi_Q}\) associated with \(\Pi_Q\), using the conventions and hypotheses of the cited Oberwolfach contribution. The target is the full algebra. The supplied record does not specify a proposed presentation, comparison algebra, or other criterion for a complete structural description. It also does not supply enough information to fix the quantification over \(Q\), the homology theory and coefficients, equivariance, grading shifts, or multiplication conventions. Consequently this is a structural research question with unresolved formulation parameters, rather than a fully specified conjectural isomorphism.

\textbf{Supplied context and qualifications.} The extracted passage states an isomorphism
\[
\mathcal L_0 A_{\Pi_Q}\cong U\!\left(\mathfrak g_{\widetilde Q,\widetilde W}\right),
\]
where \(\mathcal L_0 A_{\Pi_Q}\) is described as the subalgebra of ``perverse-exact raising operators,'' \(U\) denotes a universal enveloping algebra, and \(\mathfrak g_{\widetilde Q,\widetilde W}\) is described as the BPS Lie algebra for the tripled quiver. Here \(\widetilde Q\) is obtained by adding a loop \(\omega_i\) at each vertex of \(\overline Q\). The definitions of the subalgebra and BPS Lie algebra are absent, and the extracted formula for the potential \(\widetilde W\) is corrupted. The supplied literature assessment instead describes an identification of this subalgebra with a BPS Lie algebra; this wording does not distinguish that Lie algebra from its universal enveloping algebra and requires review.

The surrounding extract also mentions a lifted positive Heisenberg action, a lift of a BPS Lie algebra action to a Lie algebra object \(\mathfrak g_{\Pi_Q}\) in a category of perverse sheaves, an isomorphism \(H(\mathfrak g_{\Pi_Q})\cong\mathfrak g_{\widetilde Q,\widetilde W}\), and semisimplicity with a lift to a pure weight-zero mixed Hodge module. The spaces, morphisms, stability parameter, and action formulas are not adequately defined in the record. A subsequent condition concerning summands with zero-dimensional support is truncated. These are supplied contextual claims, not independently verified statements or additional hypotheses reconstructed here.
\subsection{Short English statement}
Describe the full cohomological Hall algebra \(A_{\Pi_Q}\) of a preprojective algebra. The supplied context describes a particular subalgebra through a universal enveloping algebra, while the question asks about the entire algebra. The intended quivers and algebra conventions remain unspecified. The supplied status is partially solved; that dated assessment has not been independently verified.
\subsection{Sources}
\begin{itemize}
\item Supplied local UnsolvedMath registry record, permanent ID \(30004882\), exact source problem code \texttt{catalogue entry}, title: \emph{Structure of Cohomological Hall Algebras of Preprojective Algebras}. Collection: Oberwolfach Reports Open Problems. No individual proposer is supplied; the proposed year is recorded as 2021.
\item \emph{Enveloping Algebras and Geometric Representation Theory} (2022), Oberwolfach Reports, pp.\ 2827--2891, DOI \texttt{10.4171/owr/2021/52}. Supplied report reference and literature-source link: \url{https://doi.org/10.4171/owr/2021/52}. No contribution-level page or individual author is identified in the supplied record.
\item Davison, supplied literature-source link: \url{https://arxiv.org/abs/2007.03289}. No paper title or precise theorem locator is supplied. The extracted passage contains a reference labelled [3], but its bibliography entry is absent; its identification with this link is not established by the supplied material.
\item Literature assessment checked at \texttt{2026-08-22}: the supplied curation labels the problem partially solved and says that the perverse-exact subalgebra is identified with a BPS Lie algebra while the structure of the entire algebra remains unresolved. This is supplied dated curation, not an independently verified literature assessment. The registry also labels the statement \texttt{corrected\_verified} and reports a direct check of the official report; those are supplied verification claims, not checks performed for this draft.
\end{itemize}
\end{document}
